%! Departures from Linearity — Unit 2, Lesson 2.5 — Homework
% GRADUAL RELEASE: this is the lesson's SCORED practice and the LAST component in the packet.
% Paper homework is the DEFAULT for this course; DeltaMath is an occasional override that
% replaces this page and strikes its score cell on the cover. Started in the last eight minutes
% of the period and finished at home. The AP Practice component before it stays unscored.
% THIRD CONTEXT: the notes teach on coffee shops and an antibiotic, Guided Practice runs on a
% lemonade stand, and these items are on apartment rents and phone resale values.
% ONE PAGE (compact budget): Part A (influence, 3 items incl. the AP-style MC with a required
% justification), Part B (the crux restated + a log-model prediction + evidence, 3 items), then
% the spiralbox — Lesson 2.5 closes Unit 2, so it previews the unit test and Unit 3.
% \answerspace{H}{} reserves exactly H in the blank; the key passes the answer as the second
% argument, so the two files paginate identically by construction. Keep the macro and every
% \boxguard byte-identical in the blank and the key.
\documentclass[12pt]{article}
\usepackage{apstats-article}
\usepackage{apstats-boxes}
\newcommand{\answerspace}[2]{\par\nopagebreak\noindent\begin{minipage}[t][#1][t]{\linewidth}%
  \color{keyred}\bfseries #2\end{minipage}\par}

\begin{document}

\pageheader{Unit 2, Lesson 2.5}{Homework}
% NAMESTRIP: no \namedateperiod here — the name row lives on the cover only.

\vspace{0.06in}

% Authored BYTE-IDENTICALLY in homework/ and homework_key/.
\boxguard[8]
\begin{remindbox}
\small
\textbf{This is your graded homework} --- scored, and due the first class after two study halls.
\end{remindbox}


\boxguard[12]
\begin{notesbox}{Part A --- Apartment Rents}
Ten apartments sit 0.5 to 3.0 miles from campus; apartment P is 9 miles out and rents for
\$1{,}300. For rent ($y$, dollars) against distance ($x$, miles), all 11 give
$\hat{y} = 1328 - 14.1x$, $r = -0.32$; the 10 without P give $\hat{y} = 1516 - 130.5x$,
$r = -0.99$.
\begin{enumerate}[leftmargin=1.9em, itemsep=2pt]
  \item Is P high-leverage? Is it influential? Justify both with the numbers above.
        \answerspace{1.1cm}{}
  \item \textbf{AP-style multiple choice.} A new point has an $x$-value far from the rest but lies
        right on their line. Adding it most likely (circle, then justify):\par
        \textbf{(A)} slope changes a lot \hfill \textbf{(B)} $r$ nears 0 \hfill
        \textbf{(C)} slope barely changes\par
        \textbf{(D)} it becomes an outlier \hspace{2em} \textbf{(E)} the association disappears
        \answerspace{0.7cm}{}
\end{enumerate}
\end{notesbox}

\boxguard[12]
\begin{notesbox}{Part B --- Phone Resale Values}
\begin{minipage}[c]{0.60\linewidth}\raggedright
Resale value ($y$, dollars) of 11 used phones against age ($x$, years): $\hat{y} = 518 -
95.4x$, \ $r = -0.94$. \textbf{1.}~A classmate says, ``$r$ is close to $-1$, so the line is a
good model.'' Respond, using the residual plot at the right.
\end{minipage}\hfill
\begin{minipage}[c]{0.37\linewidth}\centering
\begin{tikzpicture}[xscale=0.9, yscale=0.0062]
  \draw[linegray] (0,0) -- (6.3,0);
  \draw[->, thick] (0,-90) -- (0,140);
  \draw[thick] (0,-90) -- (6.3,-90);
  \foreach \y in {-50,0,50,100} { \node[left, font=\tiny] at (0,\y) {\y}; }
  \foreach \x in {0,2,4,6} { \node[below, font=\tiny] at (\x,-90) {\x}; }
  \foreach \x/\y in {0/122,0.5/50,1/-12,1.5/-35,2/-67,2.5/-59,3/-61,3.5/-44,4/-26,5/29,6/103}
    \fill[navy] (\x,\y) circle (0.06 and 9);
  \node[font=\tiny] at (3.1,-140) {Age (years)};
  \node[rotate=90, font=\tiny] at (-0.8,25) {Residual (\$)};
\end{tikzpicture}
\end{minipage}
\answerspace{1.0cm}{}
\begin{enumerate}[leftmargin=1.9em, itemsep=2pt, start=2]
  \item The log of each value gives $\widehat{\log y} = 2.81 - 0.19x$, $r = -0.9997$, and a
        random residual plot. Predict the value of a 3-year-old phone.
        \answerspace{0.7cm}{}
  \item Give two pieces of evidence that the log model is more appropriate than the line.
        \answerspace{0.6cm}{}
\end{enumerate}
\end{notesbox}

% Lesson 2.5 closes Unit 2, so the preview is the unit test and Unit 3, not a next lesson.
\boxguard[4]
\begin{spiralbox}
\small
\textbf{Next: the Unit 2 test} (all five lessons --- work the sample test first), \textbf{then
Unit 3, Collecting Data}: why only a well-designed experiment earns the word \emph{cause}.
\end{spiralbox}

\end{document}
